\documentclass[10pt]{article}
\usepackage[T1]{fontenc}\usepackage{lmodern,amsmath,amssymb,amsthm,microtype,booktabs}
\usepackage[margin=.88in]{geometry}\usepackage[colorlinks=true,linkcolor=black,citecolor=black,urlcolor=blue]{hyperref}
\setlength{\parindent}{0pt}\setlength{\parskip}{4pt}\setlength{\emergencystretch}{2em}
\newtheorem{theorem}{Theorem}\newtheorem{lemma}[theorem]{Lemma}\newcommand{\C}{\mathbb C}
\newcommand{\T}{\mathcal T}\DeclareMathOperator{\Tr}{Tr}\DeclareMathOperator{\Ad}{Ad}
\title{Monomial or universal: the projective groups of qutrit unitary error bases}
\author{Research manuscript}\date{20 November 2026}
\begin{document}\maketitle
\begin{abstract}\noindent A nine-element orthogonal basis of qutrit unitaries either admits a monomial representative or its projective differences generate a dense subgroup of $PU(3)$. In the monomial case the closure is finite or conjugate to the full diagonal torus extended by a cyclic coordinate permutation. The attainable finite orders are exactly $9(a^2-ab+b^2)$, for integers $a,b$ not both zero. We construct an explicit example in the dense branch using the seven translates of a cyclic difference set; its universality has an independent elementary proof. We use the finite primitive plane-group classification and a reproducible finite clique check, but no numerical approximation or classification of all qutrit unitary bases.
\end{abstract}
\section{A basis whose errors generate every projective qutrit gate densely}
A unitary error basis (UEB) $S\subset PU(3)$ is represented by nine pairwise Hilbert--Schmidt orthogonal unitaries. Individual phases and ordering do not matter; equivalence is $S\mapsto ASB$ for fixed projective unitaries. The intrinsic group $G=\langle SS^{-1}\rangle$ transforms by conjugation and is generated by any normalized representative $S s^{-1}$, $s\in S$. This distinction excludes trivial universality induced by extraneous left or right factors. Werner's shift-and-multiply construction and inequivalent qutrit bases~\cite{Werner,Nayak}, and the continuous SIC-derived wild family~\cite{S}, are prior. Here \emph{universal} means dense for one-qutrit gates modulo phase, not multi-qudit computational universality.

Let $\omega=e^{2\pi i/3}$, $\eta=e^{2\pi i/7}$, $Z=\operatorname{diag}(1,\omega,\omega^2)$ and $D=\operatorname{diag}(\eta,\eta^2,\eta^4)$. Put $u=(-13+3i\sqrt3)/14$, $v=(1,1,1)^T/\sqrt3$, $H=I+(u-1)vv^\dagger$, and take the set
\[
          S=\{I,Z\}\ \cup\ \{B_j=D^j H Z^2 D^{-j}:0\le j<7\}.
\]
\begin{theorem}This normalized UEB generates a projectively dense subgroup.
\end{theorem}
\begin{proof}The operator $H$ is unitary ($|u|=1$), with diagonal entries $(3+\omega)/7$ and off-diagonal entries $(-4+\omega)/7$. Their squared moduli are $1/7$ and $3/7$. The common diagonal of $B_j$ is proportional to $Z^2$, orthogonal to $I,Z$. The six ordered differences of $1,2,4$ exhaust the nonzero residues modulo seven. Hence $\Tr B_j^\dagger B_k=(3/7)(1+\sum_{r=1}^6\eta^{r(k-j)})=0$ for $j\ne k$.

The products $B_j Z$ are $I+(u-1)v_jv_j^\dagger$ with $v_j=D^j v$. Since $u+u^{-1}=-13/7$ is a rational noninteger, $u$ is not a root of unity. Powers of each of these products are dense in the circle of all phases on the line $\C v_j$ acting trivially on its orthogonal complement. The first three $v_j$ span by the Vandermonde determinant and are pairwise nonorthogonal: for distinct $j,k$ their squared overlap is $(3+\sum_{r=1}^6\eta^{r(k-j)})/9=2/9$.

For completeness, overlapping rank-one phase gates generate all unitaries on their span. Suppose a closed unitary group has Lie algebra containing $\mathfrak u(W)$ on a subspace $W$ (zero on $W^\perp$), and $i vv^\dagger$ with $v=w+x$, $0\ne w\in W$, $0\ne x\perp W$. Commuting once and twice with $iP_W$ isolates $wx^\dagger-xw^\dagger$ and $i(wx^\dagger+xw^\dagger)$, up to signs. Conjugating by $U(W)$ gives links to every direction in $W$, and bracketing the two links gives, up to nonzero normalization, the difference of the diagonal generators on $\C w$ and $\C x$. Hence one obtains $\mathfrak u(W+\C x)$. Starting with $\C v_0$ and adjoining $v_1,v_2$ gives $\mathfrak u(3)$ in the ordinary unitary closure of these chosen lifts, and therefore $PU(3)$ projectively.
\end{proof}
\section{Which intrinsic orders are finite?}
\begin{theorem}\label{tri}For any qutrit UEB exactly one of the following occurs.
\begin{enumerate}\setlength{\itemsep}{0pt}
\item $G$ is finite. Its basis is monomial up to equivalence, and its order belongs to $\{9(a^2-ab+b^2):a,b\in\mathbb Z,\ (a,b)\ne(0,0)\}$. Every listed order occurs.
\item $\overline G$ is unitarily conjugate to $\T\rtimes C_3$, where $\T$ is the projective diagonal torus and $C_3$ cyclically permutes coordinates. The basis is monomial up to equivalence.
\item $\overline G=PU(3)$, and the basis is not monomial up to equivalence.
\end{enumerate}
\end{theorem}
We first handle the monomial case intrinsically. Right-normalize by a member so $I$ belongs to the set. In monomial coordinates the permutations of supports must then be even: a transposition has exactly one fixed point, so its matrix cannot be traceless. Write $X$ for the cyclic shift. There must be three elements of each possible support $X^r$, since the diagonal matrices have dimension three. Any traceless diagonal unitary is projectively $Z$ or $Z^2$: three unit numbers adding to zero form an equilateral triangle. It follows that the rows of the basis are $t_r X^r K$ with $t_r\in\T$, $K=\langle Z\rangle$, and $t_0=1$. Write $R(t)=XtX^{-1}$. Conjugating by a diagonal unitary can set $t_1=1$, as $b\mapsto R(b)/b$ is onto $\T$ (rescale the desired three phases to have product one, then solve successive coordinate ratios). Thus the basis becomes
\[
                          K\ \cup\ XK\ \cup\ T X^2 K,\qquad T\in\T .
\]
Its group has index three over its diagonal part $L=\langle K,T,R(T),R^2(T)\rangle$. In $\T/K$, $L/K$ is the module generated by the point $t=TK$ under $\mathcal O=\mathbb Z[\omega]$, with $R$ acting as $\omega$: the relation $1+R+R^2=0$ holds in additive notation because the cyclic product of any diagonal is scalar.

The complex torus $\T/K$ is $\mathcal O$-equivariantly isomorphic to $\C/\mathcal O$. Indeed the real Lie algebra is the plane of zero-sum triples with a $120^\circ$ rotation. Identify it with $\C$ making this rotation multiplication by $\omega$. The period lattice, including the finite quotient by $K$, is an $\mathcal O$-stable rank-two integer lattice. Its rational span is one-dimensional over $\mathbb Q(\omega)$; after a complex scaling the lattice is a fractional $\mathcal O$-ideal, principal since $\mathcal O$ is Euclidean for the norm. After one more scaling it is $\mathcal O$. If $t$ is torsion, $\mathcal O t\simeq\mathcal O/(\alpha)$ for some $\alpha=a+b\omega\ne0$, and has $a^2-ab+b^2$ elements. Conversely choose $t=1/\alpha+\mathcal O$ and lift to $T$. This realizes every finite order in the theorem, including order nine for $\alpha$ a unit. If $t$ is nontorsion, the closure of $\mathcal O t$ is infinite and is invariant under $R$. Its identity component is a nonzero subtorus; a $120^\circ$ rotation has no invariant real line, so it is all of $\T/K$. Hence $\overline L=\T$ and case 2 results.

\paragraph{Eliminating the finite primitive groups.}
The natural $SU(3)$-lift acts irreducibly because the normalized UEB spans $M_3(\C)$. An imprimitive irreducible action in dimension three permutes three independent lines. They are orthogonal in the unitary realization: the sum of their unit-vector projectors commutes with the group, hence is scalar and equals $I$ by trace. Thus such a basis is already monomial after unitary conjugation.

For finite primitive groups we use Blichfeldt's classification~\cite{Blich}: the projective group is one of $A_5$, $\operatorname{PSL}_2(7)$, $A_6$, or a subgroup of the Hessian group $\operatorname{ASL}_2(3)=\mathbb F_3^2\rtimes\operatorname{SL}_2(3)$ (the primitive Hessian orders are 36, 72 and 216). Each nonidentity member of a normalized qutrit UEB, and each quotient of two distinct members, has projective order three by the equilateral eigenvalue observation. In each abstract group form the graph on order-three elements, with $x\sim y$ if $x^{-1}y$ has order three. Every UEB would lie in a clique, with the identity adjoined. Exact enumeration gives:
\[
\begin{array}{c|r|l|l}
 \text{ambient group}&\#\text{ vertices}&\text{maximal sizes including }1\ (\text{counts})&
 \text{generated groups for size}\ge9\\\hline
 A_5&20&3\ (10),\ 4\ (20)&\text{none}\\
 \operatorname{PSL}_2(7)&56&3\ (28),\ 4\ (784)&\text{none}\\
 A_6&80&5\ (60),\ 9\ (10)&\text{ten of order }9\\
 \operatorname{ASL}_2(3)&80&27\ (4)&\text{four of order }27
\end{array}
\]
In the last two rows every sufficiently large maximal clique including identity is actually a subgroup (size equals its generated order). Thus a UEB in an exceptional group cannot generate any primitive group on the list.

For clarity, the executable finite check uses even permutations of five/six letters, fractional linear determinant-one transformations of $\mathbb P^1(\mathbb F_7)$, and determinant-one affine transformations of $\mathbb F_3^2$ respectively. Products are composition of permutations, and vertex membership is checked by cubing. The following complete maximal-clique recursion is used. Start with $C=\varnothing$, possible vertices $P$ all vertices, excluded vertices $E=\varnothing$. If $P\cup E=\varnothing$, record $C$; otherwise choose $u\in P\cup E$. For each $v\in P\setminus N(u)$ recurse on $(C+v,P\cap N(v),E\cap N(v))$, then move $v$ from $P$ to $E$. Any still-unrecorded maximal extension must contain such a $v$ or it could be extended by $u$ (or was excluded already). Multiplying by the members from the identity until no new elements appear gives each generated order. The small check tests only necessary order-three conditions, so it does not assume order three alone implies tracelessness in every representation.

\paragraph{The infinite closed possibilities.}
Lift the compact projective closure to $SU(3)$, writing $K_0$ for the identity component. An infinite compact Lie group has positive-dimensional identity component. The full group permutes the isotypic summands for $K_0$ transitively by irreducibility. If there are multiple summands they must therefore be three lines, giving the monomial case. If the restriction is isotypic with a nontrivial multiplicity, the primeness of three means $K_0$ acts by scalars, impossible for a positive-dimensional subgroup of $SU(3)$. It remains to consider an irreducible connected compact group. Its Lie algebra is semisimple since its center acts by scalars. Its rank is at most two and its dimension at most eight. A product of two simple factors has no faithful irreducible representation of dimension three, since both nontrivial irreducibles have dimensions at least two and dimensions multiply. The simple compact types of rank at most two are $A_1,A_2,B_2,G_2$ of dimensions $3,8,10,14$. Thus it is all $\mathfrak{su}(3)$ or the three-dimensional irreducible of $\mathfrak{su}(2)$, the rotation representation (conjugate to $SO(3)$). The projective normalizer of this $SO(3)$ is itself: every automorphism is inner, and its centralizer in the unitary group is scalar by irreducibility (projective centralizing phases are trivial characters of connected semisimple $SO(3)$).

Finally a qutrit basis cannot be contained projectively in $SO(3)$. Lift its nine rotations to real unit quaternions $z_i$. Trace zero of distinct relative rotations means rotation angle $2\pi/3$, equivalently $|z_i\cdot z_j|=1/2$. The doubled Gram matrix has diagonal 2 and off-diagonal entries $\pm1$. Modulo two it equals $J-I$, of rank eight in size nine: its kernel consists of constant vectors. Its real rank would therefore be at least eight, contradicting the dimension four of quaternion space. This proves the trichotomy. The finite spectrum and the dense example describe generated groups, not a classification of arbitrary bases within the dense branch.

\begin{thebibliography}{9}\setlength{\itemsep}{0pt}
\bibitem{Blich} H.F.~Blichfeldt, \emph{Finite Collineation Groups}, University of Chicago Press (1917), Ch.~V; see also I.V.~Dolgachev, \href{https://www.maths.ed.ac.uk/cheltsov/ed/pdf/asia.pdf}{\emph{Finite subgroups of the plane Cremona group}}, Thm.~14 (the Hessian group is the affine special linear group).
\bibitem{Werner} R.F.~Werner, \emph{All teleportation and dense coding schemes}, \href{https://arxiv.org/abs/quant-ph/0003070}{arXiv:quant-ph/0003070}.
\bibitem{Nayak} A.~Nayak, H.~Yuen, \emph{Rigidity of superdense coding}, \href{https://arxiv.org/abs/2012.01672}{arXiv:2012.01672}.
\bibitem{S} E.~Svegborn, J.~Pauwels, N.~Gisin, A.~Pozas-Kerstjens, \emph{Uncountably many inequivalent maximally entangled measurements for two qutrits}, \href{https://arxiv.org/abs/2607.16396}{arXiv:2607.16396}.
\end{thebibliography}
\end{document}
